\chapter{Reinforcement Learning}
\label{chap:reinforcement_learning}

\glsresetall

This chapter introduces the fundamental concepts of \gls{rl}, providing the necessary foundation for the ideas developed throughout the rest of this manuscript.
% This chapter can be seen as a synthesis of the reference work \underline{Reinforcement Learning: An Introduction} \cite{suttonbarto1998}.
Its structure is as follows:
\begin{itemize}
    \item[]Section~\ref{sec:rl_historical_developments}: An overview of the milestones that led to the emergence of \gls{rl} as a unified field.
    \item[]Section~\ref{sec:rl_agent_paradigm}: A clarification of what is meant by an agent and how it relates to the \gls{rl} algorithm.
    \item[]Section~\ref{sec:rl_problem_formulation}: A formalization of the problem that \gls{rl} algorithms aim to solve.
    \item[]Section~\ref{sec:rl_learning}: An explanation of how policies are progressively refined through interaction with the environment to approach the optimal policy.
    \item[]Section~\ref{sec:rl_conclusion}: A summary of the key ideas and a discussion of the broader significance of \gls{rl}.
\end{itemize}

\section{Historical Developments}\label{sec:historical_developments}

Reinforcement learning is today presented as a unified field within artificial intelligence.  
However, it historically emerged from the gradual convergence of several initially independent research areas.  
Modern reinforcement learning mainly arises from the meeting of three threads: trial-and-error learning, optimal control, and temporal-difference learning.  
Each of these threads provided a different perspective on the problem of adaptation, but their combination progressively led to the theoretical framework used today.  

The fundamental intuition behind reinforcement learning is that of trial-and-error learning, historically inspired by the study of animal behavior.
The central idea is that an agent can learn without explicit supervision by interacting with its environment and gradually adapting its decisions based on observed consequences.
This view dates back to the beginning of the 20th century with the Law of Effect \cite{thorndike1911law}, according to which actions that lead to favorable outcomes tend to be repeated.
This approach strongly influenced early thinking in artificial intelligence and introduced the modern notions of an agent, interaction with the environment, and learning from a reward signal.
This idea appears in the speculation that a machine seeking to maximize a reward function could be guided by a pleasure–pain system \cite{turing1948}.
These early approaches remain essentially descriptive and do not yet provide a sufficiently general mathematical framework to rigorously formalize decision-making and its optimization.

The second thread comes from optimal control.  
In this field of mathematics, the goal is to determine a sequence of actions that maximizes a performance function over time.  
In the 1950s, Richard Bellman introduced dynamic programming \cite{bellman1957dp}, a mathematical formulation of sequential decision-making aimed at solving this type of problem.  
This approach relies on a recursive decomposition of the problem, allowing the quality of future decisions to be evaluated in order to construct an optimal policy.  

This line of work gradually led to the development of Markov Decision Processes (MDPs) as a standard framework for modelling sequential decision-making under uncertainty \cite{howard1960mdp}.  
MDPs provide a mathematical description of systems in which decisions influence future states and are associated with numerical rewards. 
A policy then describes the agent’s behavior by mapping each state to an action, or to a distribution over actions, with the objective of maximizing cumulative reward.  
The MDP framework remains today the standard mathematical formulation used to model reinforcement learning problems.

Methods derived from dynamic programming assume that the environment dynamics are known in advance.  
This assumption significantly limits their direct applicability to many real-world problems, as it requires an explicit representation of the transition model, often in the form of a complete state-transition graph.  
Such a representation quickly becomes infeasible in large-scale or continuous environments.  

The third thread is that of temporal-difference learning.  
Its central idea is to learn to evaluate states by gradually refining predictions over time using new experience \cite{sutton1988td}.  
Unlike classical dynamic programming, which relies on explicit knowledge of the environment’s transition model, temporal-difference methods enable incremental learning directly from experience.  
They thus provide an approximation of Bellman equations without requiring full access to the system dynamics.  
This family of methods constitutes a fundamental basis of modern reinforcement learning and has notably led to actor–critic architectures \cite{barto1983actorcritic} as well as many learning-based control methods.  

The fusion of these three threads became fully explicit in the late 1980s with the introduction of algorithms such as Q-learning \cite{watkins1989qlearning,watkins1992qlearning}.  
These works demonstrated that it is possible to learn complex behaviors solely through interaction with the environment.  
\section{Agent Paradigm}
\label{sec:agent_paradigm}

The \gls{rl} framework is built around the agent paradigm, in which the decision-making entity is referred to as the agent.
A standard definition is:
\begin{quote}
\emph{``An agent is anything that can be viewed as perceiving its environment through sensors and acting upon that environment through actuators.''} \cite{russell2021artificial}
\end{quote}

When this definition is adapted to the \gls{rl} setting, it can be described as follows.
The agent perceives its environment through observations.
These observations are then processed by a decision-making mechanism referred to as the policy.
The policy determines how the agent modifies its environment through the actions it selects.
The resulting action modifies the environment and, in turn, affects its subsequent dynamics.

In the context of \gls{rl}, the policy is obtained through learning.
The objective of the agent is to maximize the cumulative reward it receives over time through its interactions with the environment.
The procedure responsible for this process is the \gls{rl} algorithm. The algorithm progressively modifies the policy based on the experience collected through interaction.

In this manuscript, the \gls{rl} algorithm is considered as a process external to the agent.
This distinction is a deliberate conceptual choice.
Some other presentations of \gls{rl}, describe the agent as both the learner and the decision-making entity \cite{suttonbarto1998}.
We chose to distinguish the agent from the \gls{rl} algorithm because this provides a more general definition of an agent.

An agent can therefore exist without learning, for example when its behavior is entirely specified by a set of programmed rules.
Such an agent still perceives, processes information, and acts upon its environment, but its behavior is not the result of \gls {rl}.
\section{Objective}
\label{sec:objective}

Now that the main developments shaping reinforcement learning have been outlined, we turn to the formal definition of its objective.  
In essence, reinforcement learning aims to build a system that makes sequential decisions in order to maximize a cumulative reward signal over the long term.  
The purpose of this section is to translate this intuitive goal into a rigorous mathematical framework.

To this end, we proceed step by step:
\begin{enumerate}
    \item We define how problems are modeled using Markov Decision Processes in Subsection~\ref{subsec:markov_decision_processes}.
    \item We introduce the notion of policy, which formalizes the agent's behavior, in Subsection~\ref{subsec:policy}.
    \item We describe how the interaction between the agent and the environment is characterized through trajectories in Subsection~\ref{subsec:trajectory}.
    \item We present how policies are evaluated using value functions in Subsection~\ref{subsec:value_functions}.
    \item We define the ultimate objective of reinforcement learning: the optimal policy, in Subsection~\ref{subsec:optimal_policy}.
\end{enumerate}

\subsection{Markov Decision Processes}
\label{subsec:markov_decision_processes}

The RL framework is grounded in the agent paradigm, where the decision-making entity is referred to as the agent.
The agent is characterized by its ability to perceive its environment, process the information it receives, and, in turn, act upon the environment to influence future outcomes.
To formally specify the agent's task as well as its capacity to perceive and act, we model it as a Markov Decision Process (MDP).

Before providing a more detailed definition, it is useful to clarify the scope of the MDP problems considered in this manuscript.   
We restrict our analysis to the following configuration:  

\begin{itemize}
    \item We consider the single-agent learning setting.  
    If any other agents are present in the environment, they are treated as part of the environment dynamics.

    \item We focus on finite-horizon MDPs.  
    Each interaction sequence between the MDP and the policy has a defined beginning and a defined end.  
    Between these two points, the number of interaction steps may vary depending on the situation, but it always remains finite.

    \item We adopt the partially observable framework, where the agent does not have direct access to the true state of the system, but instead relies on limited observations that provide only partial information about the environment.  
\end{itemize}

The resulting framework can be formally described as a single-agent finite-horizon partially observable markov decision process.  
For simplicity, we will refer to this setting as an “MDP” throughout this manuscript.  

Now that the scope of the considered problems has been defined, we formally introduce the model.  
A MDP is defined by the tuple \((\mathcal{S}, \mathcal{O}, \mathcal{A}, \mathcal{R}, p)\):  
\begin{itemize}
    \item \textbf{\(\mathcal{S}\)}: The set of states, representing all possible configurations of the environment.  
    A state belonging to this set is denoted by \(s \in \mathcal{S}\).  
    Among these, we distinguish two special subsets: initial states, from which an interaction can begin, and terminal states, in which the interaction ends.  
    The set of initial states is denoted by \(\mathcal{S}^i\), with an initial state written as \(s^i \in \mathcal{S}^i\).  
    The set of terminal states is denoted by \(\mathcal{S}^t\), with a terminal state written as \(s^t \in \mathcal{S}^t\).
    
    \item \textbf{\(\mathcal{O}\)}: The set of observations that the agent can perceive from the environment.
    Unlike the true state \(s \in \mathcal{S}\), which is hidden from the agent, the observation \(o \in \mathcal{O}\) provides partial information about the current state.

    \item \textbf{\(\mathcal{A}\)}: The set of actions that the agent may execute.

    \item \textbf{\(\mathcal{R}\)}: The set of rewards that the agent can receive from the environment.

    \item \textbf{\(p\)}: The transition function, which defines the dynamics of the environment.  
    Given a state \(s\) and an action \(a\), this function defines a probability distribution over the joint outcome consisting of the next state \(s'\), the observation \(o\), and the reward \(r\).  
    We denote this function by \(p\), and it is defined as:
    \(
    p : \mathcal{S} \times \mathcal{A} \to \mathcal{P}(\mathcal{S} \times \mathcal{O} \times \mathcal{R}).
    \)
    The notation \((s', o, r) \sim p(s, a)\) means that the tuple \((s', o, r)\) is sampled from the distribution induced by \(p(\cdot \mid s, a)\).  
    The notation \(p(s', o, r \mid s, a)\) denotes the probability of observing the tuple \((s', o, r)\) under the distribution \(p(\cdot \mid s, a)\).
\end{itemize}

\subsection{Policy}
\label{subsec:policy}

The agent’s decision-making capability is modeled by a policy function.  
The policy defines the stochastic mapping between observations \(o \in \mathcal{O}\) and actions \(o \in \mathcal{A}\).  
We denote a policy function by \(\pi\), and it is defined as:
\(
\pi : \mathcal{O} \to \mathcal{P}(\mathcal{A}).
\)
The notation \(a \sim \pi(o)\) means that the action \(a\) is sampled from the distribution induced by \(\pi(\cdot \mid o)\).  
The notation \(\pi(a \mid o)\) denotes the probability of selecting action \(a\) under the distribution \(\pi(\cdot \mid o)\).

The policy thus fully characterizes the agent's decision-making capacity: given what it perceives, it determines what the agent does.
Solving an MDP precisely amounts to finding a policy \(\pi\) that maximizes the expected cumulative reward from any given state \(s\).
In other words, among all possible policies, we seek the one that yields the highest long-term return.

\subsection{Trajectory}
\label{subsec:trajectory}

Now that both the MDP and the policy have been defined, we can characterize the interaction between them.  
A sequence of consecutive interactions between an environment and a policy is called a trajectory.  
A trajectory is generated through the repeated interaction between the policy and the environment, following the sequence below:
\begin{enumerate}
    \item The environment provides an observation \(o\) to the policy \(\pi\).
    Based on this observation, the policy outputs a probability distribution over actions \(\pi(\cdot \mid o)\), from which an action is sampled:
    \(
    a \sim \pi(\cdot \mid o).
    \)

    \item The sampled action is executed in the environment. The environment then generates a transition:
    \(
    (s', o', r) \sim p(\cdot \mid s, a).
    \)
    It moves to a new state \(s'\), and returns a new observation \(o'\) and a reward \(r\).
\end{enumerate}

A trajectory \(t\) of length \(T\) can be written as:
\[
t = (s_1, o_1, a_1, r_1, s_2, o_2, a_2, r_2, \ldots, s_T, o_T, a_T, r_T).
\]

It is important to clearly distinguish between a trajectory distribution and its realizations.  
We denote by \(\tau\) the distribution over trajectories, and by \(t\) a particular trajectory.
The distribution \(\tau\) is determined by the initial state \(s\), the initial observation \(o\), the environment dynamics \(p\), and the policy \(\pi\), and is denoted by \(\tau(s, o, p, \pi)\).
The notation \(t \sim \tau(s, o, p, \pi)\) means that the trajectory \(t\) is sampled from the distribution induced by \(s\), \(o\), \(p\), and \(\pi\).
Finally, \(\tau(t \mid s, o, p, \pi)\) denotes the probability of observing the specific trajectory \(t\) under this distribution.

An important property of the trajectory distribution is its recursive structure.
Since both the policy \(\pi\) and the transition function \(p\) are known, the distribution over trajectories starting from \(s_1\) with observation \(o_1\) can be expressed in terms of the distribution over trajectories starting from the next state \(s_2\) with observation \(o_2\).
Specifically, the probability of a trajectory \(t = (s_1, o_1, a_1, r_1, s_2, o_2, a_2, r_2, \ldots, s_T, o_T, a_T, r_T)\) factorizes as:
\[
\tau(t \mid s_1, o_1, p, \pi) = \pi(a_1 \mid o_1) \; p(s_2, o_2, r_1 \mid s_1, a_1) \; \tau(t_{2:} \mid s_2, o_2, p, \pi),
\]
where \(t_{2:} = (s_2, o_2, a_2, r_2, \ldots, s_T, o_T, a_T, r_T)\) denotes the remaining trajectory from the second step onward.

This recursive decomposition mirrors the recursive form of the return \(g(t) = r_1 + \gamma \, g(t_{2:})\) and is the fundamental reason why value functions themselves satisfy recursive Bellman equations.

There exist two special types of trajectories.
First, a trajectory that terminates in a terminal state is called a terminal trajectory, and is denoted by \(t_t \sim \tau_t(s, o, p, \pi)\).
Second, a terminal trajectory that starts from an initial state and an initial observation is called an episode, and is denoted by \(t_e \sim \tau(s^i, o^i, p, \pi)\).

\subsection{Value Functions}
\label{subsec:value_functions}

To evaluate the ability of a policy \(\pi\) to collect rewards, we introduce the notion of return.  
The return of a trajectory \(g(t)\), is defined as the discounted sum of rewards collected along the trajectory:
\[
g(t) = r_1 + \gamma r_2 + \gamma^2 r_3 + \cdots + \gamma^{T-1} r_T = \sum_{i=1}^{T} \gamma^{i-1} r_i,
\]
where \(\gamma \in [0, 1]\) is the discount factor.  
This parameter controls the trade-off between immediate and future rewards: when \(\gamma = 0\), the agent cares only about the immediate reward \(r_1\); as \(\gamma\) approaches \(1\), future rewards are weighted almost equally to immediate ones, encouraging far-sighted behavior.

The return also admits a recursive form, which is central to the dynamic programming and reinforcement learning algorithms presented later:
\[
g(t) = r_1 + \gamma \, g(t_{2:}),
\]
where \(t_{2:} = (s_2, o_2, a_2, r_2, \ldots, s_T, o_T, a_T, r_T)\) denotes the trajectory starting from the second step.

From the notion of return, we can define functions that estimate the expected return of a policy given a state.  
These functions are called value functions.

The first one we introduce is the \textbf{state-value function} \(V_\pi\), which estimates the expected return when starting from a state \(s\) and following the policy \(\pi\):
\[
V_\pi(s, o) = \mathbb{E}\Big[g(t)\Big], \quad \text{with } t \sim \tau_t(s, o, p, \pi).
\]

The state-value function is particularly useful because it allows us to compare the desirability of different states under a given policy \(\pi\).  
If \(V_\pi(s) > V_\pi(s')\), this means that, according to policy \(\pi\), it is preferable for the agent to be in state \(s\) rather than in state \(s'\).

It is worth noting that the state-value function \(V_\pi\) is recursive, a property that follows directly from the recursive form of the return.  
Using \(g(t) = r_1 + \gamma \, g(t_{2:})\), we can write:
\[
\begin{aligned}
V_\pi(s, o) 
&= \mathbb{E}\Big[g(t)\Big], \quad \text{with } t \sim \tau_t(s, o p, \pi) \\[4pt]
&= \mathbb{E}\Big[r_1 + \gamma \, g(t_{2:})\Big] \\[4pt]
&= \mathbb{E}\Big[r_1 + \gamma \, V_\pi(s_2)\Big],
\end{aligned}
\]
where \(r_1\) is the reward obtained after taking the first action from state \(s\) according to \(\pi\), and \(s_2\) is the resulting next state.

If we express the expectation explicitly in terms of the distributions \(p\) and \(\pi\), we obtain the Bellman equation for \(V_\pi\):

\[
\begin{aligned}
V_\pi(s, o) 
&= \mathbb{E}\Big[r_1 + \gamma \, V_\pi(s_2)\Big], \quad \text{with } t \sim \tau_t(s, o, p, \pi) \\[4pt]
&= \sum_{a \in \mathcal{A}} \pi(a \mid o) 
   \sum_{\substack{s' \in \mathcal{S} \\ o' \in \mathcal{O} \\ r \in \mathcal{R}}} 
   p(s', o', r \mid s, a) \Big[\, r + \gamma \, V_\pi(s', o') \,\Big].
\end{aligned}
\]

Now that we have expressed the value of a state, we turn to the value of taking a specific action in a given state.
Indeed, it is useful to estimate the expected future gain of an action \(a\) in state \(s\), in order to compare different actions and determine which one is the most promising under a given policy \(\pi\).

The second value function we introduce is the \textbf{state-action value function} \(Q_\pi\), which estimates the expected return when starting from a state \(s\), taking an action \(a\), and then following the policy \(\pi\) for all subsequent steps. It is defined as:
\[
Q_\pi(s, a) = \mathbb{E}\Big[\, r + \gamma \, g(t) \,\Big], \quad \text{with } t \sim \tau_t(s', o, p, \pi), \quad (s', o, r) \sim p(\cdot \mid s, a).
\]

We can relate \(Q_\pi\) to \(V_\pi\) by recognizing that, after taking action \(a\) in state \(s\), the remaining return from the next state \(s'\) is precisely \(V_\pi(s')\). This yields:
\[
\begin{aligned}
Q_\pi(s, a) 
&= \mathbb{E}\Big[\, r + \gamma \, g(t) \,\Big], \quad \text{with } t \sim \tau_t(s', o, p, \pi), \quad (s', o, r) \sim p(\cdot \mid s, a) \\[4pt]
&= \mathbb{E}\Big[\, r + \gamma \, V_\pi(s', o) \,\Big] \\[4pt]
&= \sum_{\substack{s' \in \mathcal{S} \\ o' \in \mathcal{O} \\ r \in \mathcal{R}}} 
   p(s', o', r \mid s, a) \Big[\, r + \gamma \, V_\pi(s', o) \,\Big] \\[4pt]
\end{aligned}
\]

\subsection{Optimal Policy}
\label{subsec:optimal_policy}

Now that all the elements needed to formalize the reinforcement learning objective have been introduced, we can state it explicitly.  
The goal of a reinforcement learning agent is to find a policy that, for every observation, maximizes the expected cumulative reward obtained until a terminal state is reached.  
This policy is called the \textbf{optimal policy} and is denoted by \(\pi^*\).
Formally, \(\pi^*\) is defined as the policy that, for every observation \(o \in \mathcal{O}\), selects the action with the highest state-action value:
\[
\pi^*(a \mid o) =
\begin{cases}
1 & \text{if } a = \underset{a' \in \mathcal{A}}{\arg\max}\; Q_{\pi^*}(s, a'), \\[6pt]
0 & \text{otherwise},
\end{cases}
\]

where \(s\) is the true state underlying the observation \(o\).  

As we have just seen, both \(V_\pi\) and \(Q_\pi\) can be expressed explicitly in terms of the environment dynamics \(p\) and the policy \(\pi\).  
If \(p\) and \(\pi\) are fully known, one can theoretically compute the optimal policy by unfolding the graph of all possible trajectories starting from a given state \(s\) and observation \(o\).  
This is possible precisely because \(V_\pi\) and \(Q_\pi\) are recursive functions: their values can be propagated backward from terminal states through the entire state space.  
Solving an MDP in this way is known as \textbf{dynamic programming}.

In practice, however, most environments of interest have far too many states to allow an explicit representation of the dynamics \(p\).  
Storing and computing over the full transition function becomes intractable.

To overcome this limitation, instead of directly constructing the optimal policy, we adopt an iterative approach: starting from an initial policy, we progressively improve it so that it converges toward \(\pi^*\).  
This is the subject of the next section.
\section{Learning}
\label{sec:learning}

Now that the optimal policy \(\pi^*\) has been formally defined, we face the practical question of how to obtain it.
The recursive Bellman equations for \(V_\pi\) and \(Q_\pi\) suggest dynamic programming approaches, but these scale poorly to large state spaces.
Reinforcement learning offers an alternative: it seeks to solve the MDP iteratively, by generating a sequence of policies that progressively approach \(\pi^*\) through direct interaction with the environment.
In this section, we introduce two major families of reinforcement learning methods used today: critic-only methods and actor-critic methods.

\subsection{Critic-Only Methods}
\label{subsec:critic_only}

The central idea behind all critic-only approaches is to iteratively construct an approximation of the state-action value function \(Q_\pi\), which we denote by \(\hat{Q}_\pi\).
Although \(\hat{Q}_\pi\) is only an approximation of the true \(Q_\pi\), the policy always acts greedily with respect to this approximation:
\[
    \pi(a \mid o) =
    \begin{cases}
    1 & \text{if } a = \underset{a' \in \mathcal{A}}{\arg\max}\; \hat{Q}_\pi(s, o, a'), \\[6pt]
    0 & \text{otherwise}.
    \end{cases}
\]
It is precisely this way of modeling the policy that gives this family of methods its name: the policy is not represented explicitly as a standalone function, but simply emerges from scanning over the set of actions to maximize \(\hat{Q}_\pi\).
In other words, the policy has no independent existence of its own.

To give the reader a clear example of this family of methods, we present it in its simplest form: model-free, without temporal-dierence, replay buffer, bootstrapping, and using only \(\varepsilon\)-greedy exploration.
Naturally, many variants and extensions of this basic scheme exist, but the core idea remains the same.

The actual parametric form of \(\hat{Q}\) can be chosen freely, provided we have a rule to correct its predictions based on observed errors.
Since \(\hat{Q}\) aims to predict a scalar value in \(\mathbb{R}\), training it is a supervised regression task.
Any machine learning model capable of handling regression can therefore be used to represent \(\hat{Q}\).
We denote the set of parameters of the model by \(\theta\).
In deep reinforcement learning implementations, a neural network is used to model this function, in which case \(\theta\) represents the weights and biases of the network.

In this simple setting, the action selection rule is modified as follows.  
With probability \(1 - \varepsilon\), the policy selects the greedy action that maximizes \(\hat{Q}_\pi\), and with probability \(\varepsilon\) it selects an action uniformly at random.
The parameter \(\varepsilon\) controls the degree of exploration: if the policy always chooses the action it currently believes to be optimal, it may miss better alternatives that \(\hat{Q}_\pi\) has not yet learned to value correctly.
Exploring suboptimal actions from time to time is therefore essential to verify that \(\hat{Q}_\pi\) accurately reflects the true values, while still focusing most of the time on the most promising actions.
In more advanced implementations, \(\varepsilon\) can be adapted over the course of learning to control exploration dynamically, for instance by decreasing it gradually over time.

The learning process consists of two steps that repeat until convergence:
\begin{enumerate}
    \item \textbf{Data collection.} The current policy \(\pi\) interacts with the MDP. For a given starting state \(s\) and observation \(o\), the interaction produces a terminal trajectory \(t_t \sim \tau_t(s, o, p, \pi, \varepsilon)\), from which we compute the observed return \(g(t_t)\). For each state \(s\) along the trajectory, the associated return \(g(t_t)\) provides a target value.
    \item \textbf{Model update.} The pairs \((s, g(t_t))\) obtained in the previous step serve as supervised examples to train \(\hat{Q}_\theta\). We denote the improved approximation by \(\hat{Q}'_\theta\), which now predicts returns more accurately under \(\pi\). Since \(\hat{Q}'_\theta\) has changed, the greedy policy derived from it also changes, we denote this new policy by \(\pi'\).
    Because the policy is now \(\pi'\), the approximation \(\hat{Q}'_\theta\) must be trained to predict returns under \(\pi'\). This requires returning to step 1 to collect fresh interaction data generated by \(\pi'\). The two steps thus alternate until the policy and the value approximation stabilize.
\end{enumerate}

Even this simple algorithm already illustrates a difficulty inherent to all forms of reinforcement learning: the instability of the learning process.
Each update from \(\hat{Q}_\theta\) to \(\hat{Q}'_\theta\) causes the policy to shift from \(\pi\) to \(\pi'\).
For the learning to remain stable, this change should be as small as possible, so that \(\hat{Q}'_\theta\), trained to predict returns under \(\pi\), remains reasonably accurate under \(\pi'\).
However, the smaller the change, the slower the learning progresses.
A trade-off must therefore be struck between stability and learning speed.
Many improvements introduced later in the literature (e.g., target networks, experience replay) aim at stabilizing this learning loop.

Convergence guarantees for such algorithms have been proved in the strictly Markovian setting (fully observable MDP), where the value function is stored in a table and updated with a suitably decreasing learning rate.
In modern implementations, however, the value function is approximated by a neural network, which offers far better generalization and makes the method applicable to large-scale problems where tabular representations are unthinkable.
Although theoretical convergence proofs no longer hold in this approximate setting, deep critic-only methods (e.g., DQN) have demonstrated remarkable empirical success, making them the tool of choice for complex RL tasks.

However, critic-only approaches suffer from several limitations. The main one is that they do not scale well to large action spaces, since they must scan over all possible actions at each decision step to select the maximum.
To address this issue, another family of reinforcement learning algorithms exists: actor-critic methods.

\section{Conclusion}
\label{sec:rl_conclusion}

This chapter has provided an introduction to the fundamental principles of \gls{rl}.
Only the essential concepts required to understand the \gls{rl} framework have been presented.
% The broader field of \gls{rl} encompasses many additional methods and extensions.
Despite this limited scope, the two defining steps of \gls{rl} can already be observed in its simple learning procedure.
The first step consists of estimating the reward that can be obtained from the interaction between the agent and its environment in a given situation.
The second step uses this estimate to improve the current policy.
Learning therefore results from the repeated alternation of these two basic operations.

This simple iterative structure brings \gls{rl} conceptually close to \gls{dl}.
In both paradigms, simple operations can lead to complex behaviors when repeated at scale.
Their combination is particularly powerful.
A wide range of decision-making problems can be formulated as \gls{mdp}, allowing \gls{rl} to be applied to a broad range of tasks.
\gls{dl}, in turn, provides flexible function approximators capable of representing the value functions and policies required to solve these tasks.
Together, these two paradigms form the basis of deep \gls{rl}, which has achieved remarkable success across a wide range of challenging tasks.

From a broader perspective, the evolution of \gls{rl} has brought the field closer than ever to its initial biological inspiration.
In modern implementations, agents progressively improve their behavior through trial and error directly based on experience.
Moreover, their decision-making capabilities are represented by neural architectures that resemble their biological counterparts.
It is therefore not surprising that \gls{rl} has become one of the core paradigms for pursuing artificial general intelligence \cite{silver2021reward}.

As the capabilities of \gls{rl} continue to improve, a new challenge becomes increasingly important.
The ability to learn complex behaviors with limited human intervention can make the resulting agent's strategy difficult to understand.
To address this challenge, a parallel research direction has emerged, known as \gls{xrl}.
This field is the subject of the next chapter.
